%! Author = Güney Kıymaç
%! Date = 03-Aug-25

% Document

\section{Introduction}\label{sec:introduction}
    Macroeconomic simulation is a powerful, common, and well-studied tool widely used in policy testing/analysis, forecasting, and understanding the dynamics of economies.
With parametric equations, differential systems, and stochastic processes, the field is developed and standardized to a great extent.
However, as opposed to natural sciences, economies are not universally explained through concrete theoretical frameworks.
This has led to a unique pursuit of generalization in economic contexts, where parametric and weighted adjustments became inevitable for accurate modeling.
    \\

    MacroSim is a library that came to life through a thought experiment which started as ``How economics would have evolved if economic interactions were modeled with some individuality, not just as parts of a system?''.
This evolved into the thought of a process where variables are not inferred through mathematical transformations of core equations, but are modeled from the ground up from the components of the equations.
Challenging the current theory is not the goal of this thought and the library that resulted from it.
MacroSim was made as a toolkit for people who feel some benefit may come from exploring this approach.
    \\

    Through equation discovery aided by AR processes, MacroSim creates a simulation space through individually modeled features.
Importantly, no parametric or template-based equations are utilized beyond the use of linear autoregression.
This creates a powerful yet brittle discovery space, and it is certainly not recommended where accuracy is the primary concern.
Interesting/promising relationships uncovered by this process can be analyzed and tested in more robust and conventional environments.
